\section*{Reading list by section}
The papers behind each part of the deck, in the order they are discussed, with one line on why each is worth reading. Entries with an arXiv identifier are included in the paper archive; entries without one are listed with a URL. Items marked \emph{secondary source} are blog posts, model cards or commentary.

\subsection*{1. Overview: three vendors, six levels (slides 1-3)}
\begin{itemize}[leftmargin=1.2em,itemsep=3pt]\sloppy
\item \textbf{SemiAnalysis InferenceX (2026).} \textit{GPU specifications table (InferenceX)}. \url{https://inferencex.semianalysis.com/gpu-specs}. Analyst compilation, secondary source. Cross-vendor dense tensor peaks on one page; the sanity check for the three-vendor chart. Slide 2.
\item \textbf{Silicon Analysts (2026).} \textit{AMD vs NVIDIA AI GPU market share 2026}. \url{https://siliconanalysts.com/analysis/amd-vs-nvidia-ai-gpu-market-share-2026}. Analyst estimate, secondary source. The share figures the deck quotes as estimates only (NVIDIA about 80 percent, AMD under 10, Intel under 1). Slides 2 and 29.
\item \textbf{TrendForce (2026).} \textit{Rubin faces delays; Blackwell over 70 percent of 2026 shipments (April 2026)}. \url{https://www.trendforce.com/presscenter/news/20260408-13003.html}. Analyst press release, secondary source. The 2026 NVIDIA product mix behind the share discussion. Slides 2 and 33.
\item \textbf{Lockwood (2026).} \textit{Rubin R200 (processor garden page)}. \url{https://glennklockwood.com/garden/processors/R200}. Independent, open. A sourced compilation of every Rubin number with its date, including the NVL144 to NVL72 renaming; the provenance check for the Rubin column throughout. Slides 2, 4 and 31.
\item \textbf{Open Compute Project (2019).} \textit{OCP Accelerator Module (OAM) wiki and specifications}. \url{https://www.opencompute.org/wiki/Server/OAM}. Standards body, open. The OAM module and UBB baseboard specifications that AMD and Intel accelerators ship on; the board level of the six-level table. Slides 3 and 23.
\end{itemize}
\subsection*{2. The NVIDIA line as platforms (slides 4-6)}
\begin{itemize}[leftmargin=1.2em,itemsep=3pt]\sloppy
\item \textbf{NVIDIA (2020).} \textit{HGX A100 datasheet}. \url{https://www.nvidia.com/content/dam/en-zz/Solutions/Data-Center/HGX/pdf/nvidia-hgx-a100-datasheet.pdf}. Vendor datasheet. The 8-GPU board: 40/80 GB per GPU, 600 GB/s NVLink 3 per GPU, 4.8 TB/s aggregate through six NVSwitches; the first row of the platform table. Slides 4 and 21.
\item \textbf{NVIDIA (2023).} \textit{H100 Tensor Core GPU datasheet}. \url{https://openzeka.com/wp-content/uploads/2023/12/h100-datasheet-1.pdf}. Vendor datasheet (mirror). The shipping H100 peaks with the sparsity footnote: TF32 494/989, FP16 989/1,979, FP8 1,979/3,958 dense/sparse; 80 GB at 3.35 TB/s, 700 W. Slides 4, 6 and 18.
\item \textbf{CoreWeave documentation (2024).} \textit{NVIDIA HGX H100 (platform documentation)}. \url{https://docs.coreweave.com/coreweave-machine-learning-and-ai/nvidia-hgx-h100}. Cloud-provider documentation, secondary source. The 8-GPU HGX H100 totals (640 GB, 900 GB/s NVLink 4 per GPU) as deployed. Slide 4.
\item \textbf{NVIDIA (2023).} \textit{NVIDIA supercharges Hopper, the world's leading AI computing platform (press release, November 13, 2023)}. \url{https://nvidianews.nvidia.com/news/nvidia-supercharges-hopper-the-worlds-leading-ai-computing-platform}. Vendor press release. The H200: 141 GB of HBM3e at 4.8 TB/s on the unchanged GH100 die, available Q2 2024; the memory refresh in the platform table and the memory chart. Slides 4 and 5.
\item \textbf{Lenovo Press (LP1944) (2025).} \textit{NVIDIA H200 product guide}. \url{https://lenovopress.lenovo.com/lp1944.pdf}. Vendor documentation. The H200 dense and sparse peak table side by side, the clearest statement that the H200 compute equals the H100's. Slide 4.
\item \textbf{NVIDIA (2025).} \textit{GB200 NVL72 (product page)}. \url{https://www.nvidia.com/en-us/data-center/gb200-nvl72/}. Vendor page. The shipping rack with dense and sparse footnotes: 1,440/720 PFLOP/s FP4, 2,880 TFLOP/s FP64, 13.4 TB HBM3e at 576 TB/s, 130 TB/s NVLink; per GPU 186 GB at 8 TB/s. Slides 4, 6 and 22.
\item \textbf{The Register (2024).} \textit{Nvidia DGX GB200 NVL72: 120 kW and 1,200 W per GPU (March 2024)}. \url{https://www.theregister.com/2024/03/21/nvidia_dgx_gb200_nvk72}. Press, secondary source. The rack power of about 120 kW and 1,200 W per B200 in the GB200 superchip, figures NVIDIA's pages do not state. Slides 4 and 22.
\item \textbf{NVIDIA (2025).} \textit{GB300 NVL72 (product page)}. \url{https://www.nvidia.com/en-gb/data-center/gb300-nvl72}. Vendor page. 1,440/1,080 PFLOP/s FP4 sparse/dense, 100 TFLOP/s FP64 per rack (1.4 per GPU), 288 GB per GPU, 20 TB HBM3e, 1,400 W. Slides 4, 6, 12 and 22.
\item \textbf{Tom's Hardware (2025).} \textit{NVIDIA's newest top-tier AI supercomputers deployed for the first time at CoreWeave (July 2025)}. \url{https://www.tomshardware.com/tech-industry/artificial-intelligence/nvidias-newest-top-tier-ai-supercomputers-deployed-for-the-first-time-grace-blackwell-ultra-superchip-systems-deployed-at-coreweave}. Press, secondary source. The first GB300 NVL72 racks, July 3, 2025: the shipment date in the rack table. Slide 22.
\item \textbf{Tom's Hardware (2024).} \textit{Dell ships the industry's first GB200 NVL72 racks to CoreWeave (November 18, 2024)}. \url{https://www.tomshardware.com/tech-industry/artificial-intelligence/dell-reaches-milestone-with-industrys-first-enterprise-ready-nvidia-blackwell-poweredge-xe9712-server-racks}. Press, secondary source. The first complete GB200 NVL72 racks (Dell PowerEdge XE9712) delivered in November 2024; the first-shipment date in the rack table. Slide 22.
\item \textbf{Tom's Hardware (2026).} \textit{Nvidia launches Vera Rubin NVL72 at CES (January 2026)}. \url{https://www.tomshardware.com/pc-components/gpus/nvidia-launches-vera-rubin-nvl72-ai-supercomputer-at-ces-promises-up-to-5x-greater-inference-performance-and-10x-lower-cost-per-token-than-blackwell-coming-2h-2026}. Press, secondary source. The January 2026 numbers (288 GB HBM4 at 22 TB/s, 3.6 TB/s NVLink 6, 260 TB/s per rack) that the September 2026 revision later reduced. Slides 4 and 22.
\item \textbf{NVIDIA (2026).} \textit{Vera Rubin NVL72 (product page, revised September 2026)}. \url{https://www.nvidia.com/en-us/data-center/vera-rubin-nvl72}. Vendor page. The current Rubin specifications used throughout: 19.2 TB/s HBM4 per GPU, 3.0 TB/s NVLink 6, 216 TB/s and 1,400 TB/s per rack, 33 TFLOP/s FP64, 35/50 PFLOP/s NVFP4. Slides 4, 5, 6, 12, 22 and 31.
\item \textbf{Hardware Busters (2026).} \textit{Nvidia quietly lowers Vera Rubin HBM4 bandwidth and the HGX NVL8 loses a quarter of it (October 2026)}. \url{https://hwbusters.com/news/nvidia-quietly-lowers-vera-rubin-hbm4-bandwidth-and-the-hgx-nvl8-loses-a-quarter-of-it/}. Press, secondary source. The old-versus-new table of the September 2026 revision, dated with the Wayback Machine; why the deck gives the revised Rubin values. Slides 4 and 22.
\item \textbf{Wccftech (2026).} \textit{NVIDIA's official statement on the Vera Rubin spec change (October 29, 2026)}. \url{https://wccftech.com/nvidia-official-statement-on-vera-rubin-spec-change/}. Press, secondary source. NVIDIA's statement that the earlier figures were preliminary architectural targets and the new ones are what qualified production sustains. Slides 4 and 32.
\end{itemize}
\subsection*{3. AMD Instinct and CDNA (slides 7-12)}
\begin{itemize}[leftmargin=1.2em,itemsep=3pt]\sloppy
\item \textbf{AMD (2020).} \textit{Introducing AMD CDNA architecture (MI100 whitepaper)}. \url{https://www.amd.com/content/dam/amd/en/documents/instinct-business-docs/white-papers/amd-cdna-white-paper.pdf}. Vendor whitepaper. The CDNA 1 compute unit with matrix cores, 120 CUs, 11.5 TFLOP/s FP64, 32 GB HBM2 at 1.2 TB/s, three Infinity Fabric links at 92 GB/s; the GCN-to-CDNA split. Slides 7 and 12.
\item \textbf{AMD (2021).} \textit{AMD CDNA 2 architecture (MI200 whitepaper)}. \url{https://www.amd.com/content/dam/amd/en/documents/instinct-business-docs/white-papers/amd-cdna2-white-paper.pdf}. Vendor whitepaper. Two graphics compute dies per module, 400 GB/s in-package Infinity Fabric, packed FP32, FP64 at the FP32 rate; 47.9/95.7 TFLOP/s FP64 vector/matrix. Slides 7, 8 and 12.
\item \textbf{AMD (2023).} \textit{AMD CDNA 3 architecture whitepaper}. \url{https://www.amd.com/content/dam/amd/en/documents/instinct-tech-docs/white-papers/amd-cdna-3-white-paper.pdf}. Vendor whitepaper. The 3D chiplet package (8 XCDs on 4 IODs), 153 billion transistors, 256 MB Infinity Cache at 17.2 TB/s, 896 GB/s of peer links; MI300X and MI300A peaks. Slides 7, 9 and 12.
\item \textbf{AMD (2025).} \textit{AMD CDNA 4 architecture whitepaper}. \url{https://www.amd.com/content/dam/amd/en/documents/instinct-tech-docs/white-papers/amd-cdna-4-architecture-whitepaper.pdf}. Vendor whitepaper. The MI355X: N3P XCDs on N6 IODs, 185 billion transistors, MXFP4 and MXFP6 rates, matrix FP64 halved per CU, Infinity Fabric 7 x 144 GB/s. Slides 7, 10 and 12.
\item \textbf{AMD (2021).} \textit{Instinct MI200 series datasheet}. \url{https://www.amd.com/content/dam/amd/en/documents/instinct-tech-docs/instinct-mi200-datasheet.pdf}. Vendor datasheet. MI250X and MI250 peaks, 128 GB HBM2e at 3.2 TB/s, eight Infinity Fabric links totaling 800 GB/s, 500-560 W. Slides 6, 8 and 18.
\item \textbf{AMD (2022).} \textit{Instinct MI210 brochure}. \url{https://www.amd.com/content/dam/amd/en/documents/instinct-business-docs/product-briefs/instinct-mi210-brochure.pdf}. Vendor brochure. The single-die PCIe CDNA 2 part, for the CDNA 2 product range. Slide 8.
\item \textbf{AMD (2023).} \textit{Instinct MI300A datasheet}. \url{https://www.amd.com/content/dam/amd/en/documents/instinct-tech-docs/data-sheets/amd-instinct-mi300a-data-sheet.pdf}. Vendor datasheet. The APU: 228 CUs and 24 Zen 4 cores sharing 128 GB of HBM3, 61.3/122.6 TFLOP/s FP64, 550-760 W; El Capitan's processor. Slides 9 and 27.
\item \textbf{AMD (2023).} \textit{Instinct MI300X accelerator product brief}. \url{https://www.amd.com/content/dam/amd/en/documents/partner-hub/instinct/instinct-mi300x-accelerator-product-briefs.pdf}. Vendor product brief. The MI300X on the 8-GPU UBB 2.0 platform: 304 CUs, 192 GB at 5.3 TB/s, 750 W, the peak table. Slides 6 and 9.
\item \textbf{AMD (2024).} \textit{Instinct MI325X datasheet}. \url{https://www.amd.com/content/dam/amd/en/documents/instinct-tech-docs/product-briefs/instinct-mi325x-datasheet.pdf}. Vendor datasheet. 256 GB of HBM3e at 6 TB/s and 1,000 W on the MI300X silicon; the answer to the H200 in the memory chart. Slides 5 and 9.
\item \textbf{AMD (2025).} \textit{Instinct MI355X GPU brochure}. \url{https://www.amd.com/content/dam/amd/en/documents/instinct-tech-docs/product-briefs/amd-instinct-mi355x-gpu-brochure.pdf}. Vendor brochure. 78.6 TFLOP/s FP64 vector and matrix, 10 PFLOP/s dense FP4, 288 GB at 8 TB/s, 1,400 W liquid cooled. Slides 6, 10 and 12.
\item \textbf{AMD ROCm documentation (2024).} \textit{MI300 series microarchitecture}. \url{https://instinct.docs.amd.com/latest/_sources/gpu-arch/mi300.md}. Vendor documentation. The XCD has 40 CUs with 38 active; the FLOP per clock per CU table behind the CDNA 3 compute-unit diagram and the FP32 doubling. Slides 7 and 9.
\item \textbf{AMD ROCm documentation (2026).} \textit{MI250 microarchitecture}. \url{https://rocm.docs.amd.com/en/docs-7.14.1/_sources/reference/gpu-arch/mi250.md}. Vendor documentation. The two GCDs and the links between them; why software sees two GPUs per module. Slide 8.
\item \textbf{AMD ROCm documentation (2026).} \textit{MI100 microarchitecture}. \url{https://rocm.docs.amd.com/en/docs-7.14.1/reference/gpu-arch/mi100.html}. Vendor documentation. Eight shader engines of fifteen CUs: the MI100's 120 compute units. Slide 7.
\item \textbf{Smith and Alla (AMD) (2024).} \textit{MI300X generative AI accelerator and platform architecture (Hot Chips 36 slides)}. \url{https://www.hc2024.hotchips.org/assets/program/conference/day1/23_HC2024.AMD.MI300X.ASmith(MI300X).v1.Final.20240817.pdf}. Primary conference slides. The MI300X package, hybrid bonding, the Infinity Cache and the 8-GPU platform as AMD presented them. Slide 9.
\item \textbf{ServeTheHome (2025).} \textit{AMD dives deep on CDNA 4 architecture and MI350 at Hot Chips 2025}. \url{https://www.servethehome.com/amd-dives-deep-on-cdna-4-architecture-and-mi350-accelerator-at-hot-chips-2025/}. Press, secondary source. The Hot Chips 37 CDNA 4 talk: N3P XCDs on N6 IODs, the LDS growth to 160 KB, the FP64 decision. Slides 7 and 10.
\item \textbf{AMD (ISSCC 2024, reference page) (2024).} \textit{AMD Instinct MI300 series modular chiplet package: HPC and AI accelerator for exa-class systems (ISSCC 2024)}. \url{https://www.aihardware.mit.edu/amd-instincttm-mi300-series-modular-chiplet-package-hpc-and-ai-accelerator-for-exa-class-systems/}. Conference paper reference (title and abstract page). The ISSCC paper on the MI300 package; the primary source for the chiplet construction. Slide 9.
\item \textbf{Lam and Cozma (Chips and Cheese) (2023).} \textit{AMD's CDNA 3 compute architecture}. \url{https://chipsandcheese.com/p/amds-cdna-3-compute-architecture}. Independent analysis, secondary source. The CDNA 3 compute unit, cache hierarchy and the Infinity Cache measured and explained. Slides 7 and 9.
\item \textbf{Chips and Cheese (2025).} \textit{SC25: estimating MI430X FP64 and Discovery}. \url{https://old.chipsandcheese.com/?p=37639}. Independent analysis, secondary source. How the 288 TFLOP/s FP64 of the MI430X and the Discovery system size follow from AMD's statements. Slides 12 and 29.
\item \textbf{Phoronix (2025).} \textit{AMD previews Instinct MI400 series and Helios AI rack (June 2025)}. \url{https://www.phoronix.com/news/AMD-Instinct-MI400-Preview}. Press, secondary source. The June 12, 2025 numbers for MI400 and Helios (432 GB HBM4, 19.6 TB/s, 40 PFLOP/s FP4) before the 2026 launch revised them. Slide 11.
\item \textbf{Tom's Hardware (2025).} \textit{AMD debuts Helios rack-scale platform at OCP Global Summit 2025}. \url{https://www.tomshardware.com/tech-industry/amd-debuts-helios-rack-scale-ai-hardware-platform-at-ocp-global-summit-2025-promises-easier-serviceability-and-50-percent-more-memory-than-nvidias-vera-rubin}. Press, secondary source. Helios on the OCP Open Rack Wide form factor, double-wide for serviceability, and the memory comparison with Vera Rubin. Slides 11, 22 and 23.
\item \textbf{Tom's Hardware (2026).} \textit{AMD takes the wraps off Instinct MI455X (CDNA 5) and the Helios rack (July 2026)}. \url{https://www.tomshardware.com/pc-components/gpus/amd-takes-the-wraps-off-its-instinct-mi455x-ai-accelerator-cdna-5-and-helios-rack-scale-architecture-combine-to-take-the-fight-to-nvidia-in-the-data-center}. Press, secondary source. The MI455X launch of July 23, 2026: N2 XCDs, 320 billion transistors, 192 MB L2 with no Infinity Cache, 23.3 TB/s HBM4, 256 work-group processors. Slides 7, 11 and 31.
\item \textbf{StorageReview (2026).} \textit{AMD MI455X and Helios: 432 GB HBM4, 72-GPU racks (July 2026)}. \url{https://www.storagereview.com/news/amd-mi455x-and-helios-432gb-hbm4-72-gpu-racks-and-a-real-answer-to-vera-rubin}. Press, secondary source. UALink over Ethernet at 3.6 TB/s per GPU, 260 TB/s per rack, 225-245 kW; the Helios figures in the rack table. Slides 11, 21 and 22.
\item \textbf{Supermicro (2026).} \textit{AMD Helios rack datasheet}. \url{https://www.supermicro.com/datasheet/datasheet_AMD_Helios.pdf}. Vendor (integrator) datasheet. The rack bill of materials: 72 MI455X, 18 Venice CPUs, 18 Salina DPUs, Vulcano NICs, 31 TB HBM4 (listed at 1.4 PB/s). Slides 11 and 22.
\item \textbf{ServeTheHome (2026).} \textit{AMD Advancing AI 2026 keynote live coverage}. \url{https://www.servethehome.com/?p=102176}. Press, secondary source. The MI430X at 288 TFLOP/s FP64 and the MI500 (2027) and MI600 (2028) roadmap as announced. Slides 11, 12 and 29.
\item \textbf{Tech Wire Asia (2026).} \textit{AMD Advancing AI 2026: Helios, OpenAI, Meta, Anthropic}. \url{https://techwireasia.com/2026/07/amd-advancing-ai-2026-helios-openai-meta-anthropic/}. Press, secondary source. The customer commitments announced with Helios (OpenAI, Meta, Anthropic, Oracle). Slides 11 and 29.
\item \textbf{AMD (2025).} \textit{Instinct MI430X: powering the next wave of AI and scientific discovery}. \url{https://www.amd.com/en/blogs/2025/amd-instinct-mi430x-powering-the-next-wave-of-ai.html}. Vendor blog. The HPC SKU of the MI400 family, its hardware FP64 and the systems that chose it (Discovery, Alice Recoque). Slides 12 and 29.
\item \textbf{AMD (SEC Form 8-K) (2026).} \textit{AMD reports third quarter 2026 financial results}. \url{https://www.sec.gov/Archives/edgar/data/0000002488/000000248826000188/q32026991.htm}. Primary filing. Helios shipping, first Helios revenue, data-center segment of 8.5 billion dollars and 97 percent growth. Slides 11 and 29.
\item \textbf{AMD Q3 2026 earnings call (transcript) (2026).} \textit{AMD Q3 2026 earnings call transcript}. \url{https://singjupost.com/amd-advanced-micro-devices-q3-2026-earnings-call-transcript/}. Transcript, secondary source. The statement that the MI455X is shipping in volume and the customers named. Slide 11.
\item \textbf{StorageNewsletter (2026).} \textit{AMD fiscal 4Q25 and FY25 financial results}. \url{https://www.storagenewsletter.com/2026/02/18/amd-fiscal-4q25-and-fy25-financial-results/}. Press, secondary source. Data-center revenue of 16.6 billion dollars in fiscal 2025, the base for the 2026 growth figures. Slide 29.
\item \textbf{AMD investor relations (2025).} \textit{AMD and OpenAI announce strategic partnership to deploy 6 gigawatts of AMD GPUs (October 2025)}. \url{https://ir.amd.com/news-events/press-releases/detail/1260/amd-and-openai-announce-strategic-partnership-to-deploy-6-gigawatts-of-amd-gpus}. Vendor press release. The 6 GW commitment from the second half of 2026 and the warrant for up to 160 million shares. Slides 3, 11 and 29.
\end{itemize}
\subsection*{4. Intel: the Gen families, Xe and Gaudi (slides 13-16)}
\begin{itemize}[leftmargin=1.2em,itemsep=3pt]\sloppy
\item \textbf{Junkins (Intel) (2015).} \textit{The compute architecture of Intel processor graphics Gen9}. \url{https://www.intel.com/content/dam/develop/external/us/en/documents/the-compute-architecture-of-intel-processor-graphics-gen9-v1d0-843410.pdf}. Vendor whitepaper. The execution unit (7 threads, two SIMD-4 FPUs, 16 FLOP per clock), the subslice of 8 EUs with sampler and 64 KB shared local memory, the slice with 768 KB L3; the Gen9 diagram and table. Slide 13.
\item \textbf{Intel (2019).} \textit{Intel processor graphics Gen11 architecture}. \url{https://www.intel.com/content/dam/develop/external/us/en/documents/the-architecture-of-intel-processor-graphics-gen11-r1new.pdf}. Vendor whitepaper. The 64-EU Gen11 GPU of Ice Lake, 1,024 FLOP per clock, and the changes to the slice. Slide 13.
\item \textbf{Tom's Hardware (2020).} \textit{Intel drops Xe LP graphics specifications (August 2020)}. \url{https://www.tomshardware.com/news/intel-xe-lp-graphics-specs}. Press, secondary source. Xe-LP's 96 EUs and 768 ALUs in Tiger Lake and the DG1, the last Gen-style design and the first Xe. Slide 13.
\item \textbf{Intel (2022).} \textit{Introduction to the Xe-HPG architecture}. \url{https://www.intel.com/content/www/us/en/developer/articles/technical/introduction-to-the-xe-hpg-architecture.html}. Vendor documentation. The Xe-core, Xe Vector Engine and XMX naming and the Xe-HPG core (16 XVE, 16 XMX, ray-tracing unit) in the two-core diagram. Slide 14.
\item \textbf{Intel (2024).} \textit{oneAPI GPU optimization guide: Xe architecture}. \url{https://www.intel.com/content/www/us/en/docs/oneapi/optimization-guide-gpu/2024-1/xe-arch.html}. Vendor documentation. The Xe-HPC hierarchy (XVE, Xe-core, slice, stack) with lane counts, the 512 KB L1 and SLM per core, and the cross-vendor level table. Slides 14 and 18.
\item \textbf{Tom's Hardware (2022).} \textit{Intel reveals specifications for Arc Alchemist desktop GPUs (September 2022)}. \url{https://www.tomshardware.com/news/intel-reveals-specifications-for-arc-alchemist-desktop-gpus}. Press, secondary source. The ACM-G10 table: 32 Xe-cores, 21.7 billion transistors on 406 square millimeters of TSMC N6, 17.2 TFLOP/s, 16 GB at 560 GB/s, 225 W. Slide 15.
\item \textbf{Tom's Hardware (2024).} \textit{Intel Arc B580 review (December 2024)}. \url{https://www.tomshardware.com/pc-components/gpus/intel-arc-b580-review-the-new-usd249-gpu-champion-has-arrived}. Press review, secondary source. BMG-G21 on TSMC N5: 20 Xe2 cores with SIMD16 XVEs, 19.6 billion transistors on 272 square millimeters, 12 GB at 456 GB/s, 190 W. Slide 15.
\item \textbf{Tom's Hardware (2025).} \textit{Intel launches Arc Pro B50 and B60 and Project Battlematrix workstations (May 2025)}. \url{https://www.tomshardware.com/pc-components/gpus/intel-launches-usd299-arc-pro-b50-with-16gb-of-memory-project-battlematrix-workstations-with-24gb-arc-pro-b60-gpus}. Press, secondary source. The workstation Xe2 parts with 16 and 24 GB and the eight-B60 Battlematrix box for local inference. Slide 15.
\item \textbf{HotHardware (2026).} \textit{Intel CES 2026: Panther Lake (Core Ultra series 3) is a go}. \url{https://hothardware.com/news/intel-ces-2026-panther-lake-is-a-go}. Press, secondary source. The first Xe3 GPU: 12 Xe3 cores on Intel 18A, about 120 TOPS on the GPU, in a laptop processor. Slide 15.
\item \textbf{Kaplan (Intel) (2024).} \textit{Gaudi 3 (Hot Chips 2024 slides)}. \url{https://www.hc2024.hotchips.org/assets/program/conference/day1/60_HC2024.Intel.RomanKaplan.Gaudi3-0826.pdf}. Primary conference slides. Two dies, 8 matrix engines and 64 tensor processor cores, 128 GB HBM2e at 3.67 TB/s, 24 x 200 GbE on the die; the Gaudi 3 diagram. Slide 16.
\item \textbf{The Register (2022).} \textit{Intel's Habana unit reveals Gaudi 2 (May 2022)}. \url{https://www.theregister.com/2022/05/10/intels_habana_unit_reveals_new/}. Press, secondary source. Gaudi 2: 7 nm, 24 TPCs, 96 GB HBM2e at 2.45 TB/s, 24 x 100 GbE, 600 W. Slide 16.
\item \textbf{Tom's Hardware (2024).} \textit{Intel says it will miss its 500 million dollar Gaudi goal (October 2024)}. \url{https://www.tomshardware.com/tech-industry/artificial-intelligence/intel-says-it-will-miss-its-ai-goals-with-gaudi-3-unbaked-software-leaves-intels-usd500-million-ai-goal-unachievable-as-competitors-rake-in-billions}. Press, secondary source. The withdrawn 2024 revenue target and the software reasons Intel gave; the market outcome of the Gaudi line. Slides 16 and 20.
\end{itemize}
\subsection*{5. Xe-HPC: Ponte Vecchio, Aurora and after (slides 17-20)}
\begin{itemize}[leftmargin=1.2em,itemsep=3pt]\sloppy
\item \textbf{Intel (2023).} \textit{Intel Data Center GPU Max 1550 specifications (ARK)}. \url{https://www.intel.com/content/www/us/en/products/sku/232873/intel-data-center-gpu-max-1550/specifications.html}. Vendor specification page. 128 Xe-cores, 128 GB HBM2e at 3,276.8 GB/s, 600 W OAM, launch date January 2023. Slides 17, 18 and 31.
\item \textbf{Intel Labs (2022).} \textit{Intel at ISSCC 2022: Ponte Vecchio (Gomes et al.) and other papers}. \url{https://intel.com/content/www/us/en/research/blogs/isscc-2022.html}. Vendor research blog. The ISSCC 2022 paper on Ponte Vecchio as a multi-tile 3D stacked processor; the primary source for the tile census. Slide 17.
\item \textbf{TechPowerUp (2022).} \textit{Intel details Ponte Vecchio: 63 tiles, 600 W TDP and lots of bandwidth (ISSCC 2022)}. \url{https://techpowerup.com/292250/intel-details-ponte-vecchio-accelerator-63-tiles-600-watt-tdp-and-lots-of-bandwidth}. Press, secondary source. The tile census (47 active, 63 with thermal tiles), the five process nodes and the 408 MB of L2 as presented at ISSCC. Slide 17.
\item \textbf{The Register (2022).} \textit{Intel Ponte Vecchio at Hot Chips 34 (August 2022)}. \url{https://www.theregister.com/2022/08/23/intel_ponte_vecchio/}. Press, secondary source. 52 TFLOP/s FP64 equal to FP32 and 419 TFLOP/s TF32 on the XMX engines, the figures the deck uses for the Max 1550. Slides 17 and 18.
\item \textbf{Tom's Hardware (2021).} \textit{Intel Ponte Vecchio and Xe-HPC architecture: built for big data (Architecture Day 2021)}. \url{https://www.tomshardware.com/features/intel-ponte-vecchio-and-xe-hpc-architecture-built-for-big-data}. Press, secondary source. The stack composition (8 compute tiles, 4 RAMBO cache tiles, a base tile, Xe Link tiles) and the Xe-HPC core as first described. Slides 14 and 17.
\item \textbf{TechPowerUp GPU database (2023).} \textit{Intel Data Center GPU Max 1550}. \url{https://www.techpowerup.com/gpu-specs/data-center-gpu-max-1550.c4068}. Database, secondary source. The 52.43 TFLOP/s FP32 and FP64 figures, clock and lane counts behind the derived arithmetic. Slide 18.
\item \textbf{Lockwood (2025).} \textit{Ponte Vecchio (processor garden page)}. \url{https://glennklockwood.com/garden/processors/PVC}. Independent, open. The Xe-HPC hierarchy with sources and measured DGEMM on Aurora; the provenance check for the PVC column. Slides 17, 18 and 31.
\item \textbf{Phoronix (2022).} \textit{Intel introduces Xeon Max and Data Center GPU Max (November 2022)}. \url{https://www.phoronix.com/review/intel-xeon-max-dc/3}. Press, secondary source. The Max series SKUs as announced (1100, 1350, 1550) and the Xeon Max with HBM used in Aurora. Slides 17 and 19.
\item \textbf{Argonne Leadership Computing Facility (2026).} \textit{Aurora system overview (machine overview)}. \url{https://docs.alcf.anl.gov/aurora/hardware-overview/machine-overview}. Facility documentation. The node (two Xeon Max 9470C, six Max 1550, eight Slingshot NICs), 10,624 nodes in 166 racks, 63,744 GPUs. Slides 19 and 27.
\item \textbf{Goto et al. (2026).} \textit{HPL and HPL-MxP at exascale on Aurora}. arXiv:2604.09517. The run campaigns behind Aurora's TOP500 entries: 0.585 EFLOP/s on 5,439 nodes, 1.012 on 9,234, and 11.64 EFLOP/s HPL-MxP. Slide 19.
\item \textbf{HPCwire (2023).} \textit{Aurora the survivor: exascale supercomputer arrives after eight years (November 2023)}. \url{https://www.hpcwire.com/2023/11/13/aurora-the-survivor-exascale-supercomputer-arrives-after-eight-years-of-doom/}. Press, secondary source. The 2015 to 2023 history: Knights Hill, the 2017 re-plan, the 7 nm slip, installation in 2023. Slide 19.
\item \textbf{Intel (2023).} \textit{Intel systems in the November 2023 TOP500 (newsroom)}. \url{https://www.intel.com/content/www/us/en/newsroom/news/top500-supercomputers-nov-2023.html}. Vendor newsroom. Aurora's first entry and the other PVC systems (Dawn at Cambridge, SuperMUC-NG Phase 2 at LRZ). Slide 19.
\item \textbf{Argonne Leadership Computing Facility (2025).} \textit{Argonne researchers and supercomputers contribute to four Gordon Bell prize finalist projects (2025)}. \url{https://alcf.anl.gov/news/argonne-researchers-and-supercomputers-contribute-four-projects-named-gordon-bell-prize}. Facility news. The four 2025 Gordon Bell finalists run on Aurora, the application-scale evidence for the machine. Slide 19.
\item \textbf{Hatanpaa et al. (2025).} \textit{AERIS: Argonne Earth Systems Model for reliable and skillful predictions}. arXiv:2509.13523. The weather and climate model that sustained 10.21 EFLOP/s in mixed precision on 10,080 Aurora nodes. Slide 19.
\item \textbf{insideHPC (2023).} \textit{Intel alters HPC-AI roadmap: Rialto Bridge GPU discontinued (March 2023)}. \url{https://insidehpc.com/2023/03/intel-alters-hpc-ai-roadmap-rialto-bridge-gpu-discontinued/}. Press, secondary source. The cancellation of Ponte Vecchio's successor and the roadmap reset toward Falcon Shores. Slides 3 and 20.
\item \textbf{Tom's Hardware (2025).} \textit{Intel cancels Falcon Shores GPU for AI workloads; Jaguar Shores to be successor (January 2025)}. \url{https://www.tomshardware.com/tech-industry/artificial-intelligence/intel-cancels-falcon-shores-gpu-for-ai-workloads-jaguar-shores-to-be-successor}. Press, secondary source. The January 30, 2025 decision to keep Falcon Shores as an internal test chip and the naming of Jaguar Shores. Slides 3 and 20.
\item \textbf{Phoronix (2025).} \textit{Intel announces Crescent Island (October 2025)}. \url{https://www.phoronix.com/review/intel-crescent-island}. Press, secondary source. The Xe3P inference GPU with 160 GB of LPDDR5X announced at OCP 2025. Slides 15 and 20.
\item \textbf{StorageReview (2026).} \textit{Intel GPU D-series: 480 GB LPDDR5X, 1.3 PFLOP/s FP8, 350 W PCIe card (October 2026)}. \url{https://www.storagereview.com/news/intel-gpu-d-series-480gb-lpddr5x-1-3-pflops-fp8-350w-pcie-card}. Press, secondary source. The GPU D-series as named and specified in October 2026 (32 Xe3P cores, 1.5 TB/s, sampling); its current status. Slides 20 and 32.
\item \textbf{Intel (2026).} \textit{Intel reports third-quarter 2026 financial results}. \url{https://www.intc.com/news-events/press-releases/detail/1786/intel-reports-third-quarter-2026-financial-results}. Primary filing (press release). Data center and AI revenue of 7.3 billion dollars, growth on CPUs, the D-series as the only accelerator mentioned. Slides 16 and 20.
\item \textbf{Intel (2025).} \textit{NVIDIA and Intel to develop AI infrastructure and personal computing products (September 2025)}. \url{https://www.intc.com/news-events/press-releases/detail/1750/}. Vendor press release. The 5 billion dollar NVIDIA investment and the plan for x86 CPUs with NVLink. Slides 20, 23 and 29.
\end{itemize}
\subsection*{6. Fabrics, racks and standards (slides 21-23)}
\begin{itemize}[leftmargin=1.2em,itemsep=3pt]\sloppy
\item \textbf{UALink Consortium (2025).} \textit{UALink 1.0 white paper (April 2025)}. \url{https://ualinkconsortium.org/wp-content/uploads/2025/04/UALink-1.0-White_Paper_FINAL.pdf}. Standards body white paper. The 200G 1.0 specification: 200 Gb/s per lane, up to 1,024 accelerators, memory-semantic operations; the open scale-up fabric. Slides 21 and 23.
\item \textbf{Ultra Ethernet Consortium (Linux Foundation) (2025).} \textit{UEC launches specification 1.0 (June 2025)}. \url{https://www.linuxfoundation.org/press/uec-launches-spec-1.0}. Standards body press release. Ultra Ethernet 1.0, the scale-out counterpart of UALink and the alternative to InfiniBand. Slide 23.
\item \textbf{NVIDIA (2025).} \textit{NVIDIA unveils NVLink Fusion (press release, May 2025)}. \url{https://investor.nvidia.com/news/press-release-details/2025/NVIDIA-Unveils-NVLink-Fusion-for-Industry-to-Build-Semi-Custom-AI-Infrastructure-With-NVIDIA-Partner-Ecosystem/default.aspx}. Vendor press release. NVLink opened to partners' CPUs and accelerators on NVIDIA's terms; the answer to UALink in the standards timeline. Slide 23.
\item \textbf{Abraham (OLCF) (2023).} \textit{Frontier architecture overview}. \url{https://www.olcf.ornl.gov/wp-content/uploads/Frontier-Architecture-Overview_Abraham.pdf}. Facility presentation. The Frontier node: one Trento CPU, four MI250X, coherent Infinity Fabric CPU-GPU links, Slingshot-11 NICs; the node diagram. Slides 8, 21 and 27.
\item \textbf{Lawrence Livermore National Laboratory (2026).} \textit{El Capitan system page}. \url{https://hpc.llnl.gov/node/1269}. Facility documentation. Node counts, APU counts and the 2,889 PFLOP/s peak of the MI300A system. Slides 9 and 27.
\item \textbf{Morgan (The Next Platform) (2024).} \textit{El Capitan blazes the trail for converged CPU-GPU compute (November 2024)}. \url{https://www.nextplatform.com/2024/11/18/el-capitan-supercomputer-blazes-the-trail-for-converged-cpu-gpu-compute/}. Press, secondary source. 11,136 nodes and 44,544 APUs at the first TOP500 entry, and the case for the APU node. Slides 9 and 27.
\end{itemize}
\subsection*{7. Software stacks and portability (slides 24-26)}
\begin{itemize}[leftmargin=1.2em,itemsep=3pt]\sloppy
\item \textbf{AMD (2026).} \textit{ROCm release history (versions page)}. \url{https://rocm.docs.amd.com/en/latest/release/versions.html}. Vendor documentation. Dated ROCm releases from 5.0 to 10.x, the source of the dates in the three-stacks panel. Slide 24.
\item \textbf{AMD (2026).} \textit{HIPIFY documentation}. \url{https://rocm.docs.amd.com/projects/HIPIFY/en/latest}. Vendor documentation. hipify-clang and hipify-perl, the CUDA-to-HIP source translators behind the porting-path discussion. Slide 24.
\item \textbf{Khronos Group (2021).} \textit{Khronos releases SYCL 2020 final specification (February 2021)}. \url{https://www.khronos.org/news/press/khronos-releases-sycl-2020-final-specification}. Standards body press release. SYCL 2020 with unified shared memory and generic back ends; the specification Intel's DPC++ implements. Slide 24.
\item \textbf{Phoronix (2020).} \textit{Intel oneAPI 1.0 officially released (September 2020)}. \url{https://www.phoronix.com/review/intel-oneapi-10}. Press, secondary source. The oneAPI 1.0 specification release and the first toolkits that followed in December 2020. Slides 13 and 24.
\item \textbf{TechCrunch (2023).} \textit{The Unified Acceleration Foundation wants an open standard for accelerator programming (September 2023)}. \url{https://techcrunch.com/2023/09/19/the-unified-acceleration-foundation-wants-to-create-an-open-standard-for-accelerator-programming/}. Press, secondary source. oneAPI's move into the Linux Foundation's UXL, the governance point in the three-stacks panel. Slide 24.
\item \textbf{Codeplay (2022).} \textit{Codeplay announces oneAPI for NVIDIA and AMD GPU hardware (December 2022)}. \url{https://codeplay.com/portal/press-releases/2022/12/16/codeplay-announces-oneapi-for-nvidia-and-amd-gpu-hardware}. Vendor press release. The SYCL plugins that run DPC++ code on CUDA and HIP back ends; the portability-layer claim. Slides 24 and 25.
\item \textbf{Intel (2026).} \textit{Migrate from CUDA to C++ with SYCL (SYCLomatic and the DPC++ compatibility tool)}. \url{https://www.intel.com/content/www/us/en/developer/tools/oneapi/training/migrate-from-cuda-to-cpp-with-sycl.html}. Vendor documentation. The CUDA-to-SYCL migration tool and Intel's claim that 80 to 90 percent of code migrates automatically. Slide 24.
\item \textbf{Reguly (2023).} \textit{Evaluating the performance portability of SYCL across CPUs and GPUs from different vendors}. arXiv:2309.10075. SC-W 2023 (P3HPC). SYCL against native on Max 1100, A100 and MI250X: DPC++ within 1.2 percent of CUDA on the A100, 15.9 percent behind HIP on the MI250X, weaker on CPUs. Slide 25.
\item \textbf{Davis et al. (2025).} \textit{Taking GPU programming models to task for performance portability}. arXiv:2402.08950. ICS 2025. Seven programming models and five proxy applications across five GPUs: portability layers reach most of native speed and the spread sits in compilers and runtimes. Slide 25.
\item \textbf{ORNL and SNL (ORNL/TM-2025/4398) (2025).} \textit{Experiences with SYCL on AMD GPUs with Kokkos}. \url{https://www.osti.gov/biblio/3016989}. Laboratory technical report. Kokkos with DPC++ and AdaptiveCpp back ends on AMD GPUs; DPC++ less mature on AMD than on Intel, HIP the reference. Slide 25.
\item \textbf{Patel, Nishball and Knuhtsen (SemiAnalysis) (2024).} \textit{MI300X vs H100 vs H200 benchmark part 1: training (December 2024)}. \url{https://newsletter.semianalysis.com/p/mi300x-vs-h100-vs-h200-benchmark-part-1-training}. Analyst report, secondary source. Achieved BF16 GEMM against peak (about 720 of 989 TFLOP/s on the H100, about 620 of 1,307 on the MI300X), 64 against 450 GB/s between GPUs; the software gap in numbers. Slides 21 and 26.
\end{itemize}
\subsection*{8. Systems and procurements (slides 27-29)}
\begin{itemize}[leftmargin=1.2em,itemsep=3pt]\sloppy
\item \textbf{TOP500 (2026).} \textit{TOP500 list, June 2026}. \url{https://top500.org/lists/2026/06}. Benchmark list, open. The June 2026 top ten with accelerators, Rmax, Rpeak and power; the exascale and TOP500 tables. Slides 2, 12, 27 and 28.
\item \textbf{TOP500 (2026).} \textit{HPCG results, June 2026}. \url{https://top500.org/lists/hpcg/2026/06/}. Benchmark list, open. The HPCG ranks that complement HPL for the same machines. Slide 28.
\item \textbf{TOP500 (2026).} \textit{Green500 list, June 2026}. \url{https://top500.org/lists/green500/green500-june-2026/}. Benchmark list, open. GFLOP/s per watt: the Grace Hopper systems at the top and the MI300A systems near 66. Slide 28.
\item \textbf{HPE (2026).} \textit{HPE delivers six of the ten most powerful supercomputers (June 2026)}. \url{https://www.hpe.com/us/en/newsroom/blog-post/2026/06/hpe-delivers-six-out-of-ten-of-worlds-most-powerful-supercomputers.html}. Vendor blog. The HPE Cray EX share of the top ten and the HPL-MxP order (El Capitan, Frontier, Aurora). Slides 27 and 28.
\item \textbf{CSC (LUMI) (2021).} \textit{LUMI's full system architecture revealed}. \url{https://www.lumi-supercomputer.eu/lumis-full-system-architecture-revealed/}. Facility news. 2,978 GPU nodes of four MI250X each; the European Frontier sibling. Slides 8 and 28.
\item \textbf{CINES (2024).} \textit{Adastra's architecture}. \url{https://dci.dci-gitlab.cines.fr/webextranet/architecture/index.html}. Facility documentation. MI250X and MI300A partitions in one system; another AMD HPC deployment. Slide 8.
\item \textbf{HLRS (2025).} \textit{Hunter}. \url{https://www.hlrs.de/solutions/systems/hunter}. Facility page. 752 MI300A, the transition system before Herder. Slides 8 and 29.
\item \textbf{Oak Ridge National Laboratory (2025).} \textit{ORNL, AMD and HPE to deliver DOE's newest AI supercomputers, Discovery and Lux}. \url{https://www.ornl.gov/news/ornl-amd-and-hpe-deliver-does-newest-ai-supercomputers-discovery-and-lux}. Laboratory press release. Lux (MI355X, early 2026) and Discovery (MI430X with Venice on HPE Cray GX5000, 2028). Slides 10, 12 and 29.
\item \textbf{AMD and Eviden (2025).} \textit{AMD and Eviden to power Europe's new exascale supercomputer, the first in France}. \url{https://newsroom.amd.com/news/amd-and-eviden-to-power-europe-s-new-exascale-supe/}. Vendor press release. Alice Recoque: MI430X with Venice CPUs in 94 racks on the BXI network. Slides 12 and 29.
\end{itemize}
\subsection*{9. Benchmarks (slide 30)}
\begin{itemize}[leftmargin=1.2em,itemsep=3pt]\sloppy
\item \textbf{AMD (2024).} \textit{Engineering insights: unveiling MLPerf results on AMD Instinct MI300X}. \url{https://www.amd.com/en/blogs/2024/engineering-insights-unveiling-mlperf-results-on.html}. Vendor blog. AMD's first MLPerf submission (Inference 4.1, August 2024): eight MI300X within 2 to 3 percent of a DGX H100 on Llama 2 70B. Slide 30.
\item \textbf{AMD ROCm blog (2025).} \textit{MI325X accelerates MLPerf Inference v5.0}. \url{https://rocm.blogs.amd.com/artificial-intelligence/mi325x-accelerates-mlperf-inference/README.html}. Vendor blog. The first MI325X results, April 2025. Slide 30.
\item \textbf{AMD ROCm blog (2025).} \textit{Technical dive into AMD's MLPerf Inference v5.1 submission}. \url{https://rocm.blogs.amd.com/artificial-intelligence/mlperf-inference-v5.1/README.html}. Vendor blog. The first MI355X and first MXFP4 results (September 2025), Llama 2 70B at 2.7 times the MI325X. Slides 10 and 30.
\item \textbf{AMD (2026).} \textit{AMD delivers breakthrough MLPerf Inference 6.0 results: MI355X past one million tokens per second}. \url{https://amd.com/en/blogs/2026/amd-delivers-breakthrough-mlperf-inference-6-0-results.html}. Vendor blog. The April 2026 ratios against the B200 (1.00, 0.97, 1.19) and B300 (0.92, 0.93, 1.04) on Llama 2 70B, and the 87-GPU million-token run. Slides 10 and 30.
\item \textbf{MLCommons (2025).} \textit{MLPerf Training v5.0 results (June 2025)}. \url{https://mlcommons.org/2025/06/mlperf-training-v5-0-results}. Benchmark consortium, open. AMD's first training submission and the GB200's debut in the same round. Slide 30.
\item \textbf{AMD (2025).} \textit{Accelerating AI training: MI350 series in MLPerf Training v5.1}. \url{https://www.amd.com/en/blogs/2025/accelerating-ai-training.html}. Vendor blog. Llama 2 70B LoRA in 10.18 minutes against 9.85 (B200) and 9.59 (B300), and Llama 3.1 8B at 99.7 against 93.7 minutes. Slide 30.
\item \textbf{NVIDIA (2026).} \textit{NVIDIA Blackwell tops MLPerf Training 6.0 (developer blog, June 2026)}. \url{https://developer.nvidia.com/blog/nvidia-blackwell-tops-mlperf-training-6-0-with-industry-leading-scale-and-performance/}. Vendor blog. The 8,192-GPU GB300 results (DeepSeek-V3 671B in 2.02 minutes) that set the scale of Training 6.0, with no AMD result at that size. Slide 30.
\end{itemize}
\subsection*{10. Wrap-up and the record (slides 31-33)}
\begin{itemize}[leftmargin=1.2em,itemsep=3pt]\sloppy
\item \textbf{Lockwood (2026).} \textit{GTC 2026 recap (blog, March 2026)}. \url{https://blog.glennklockwood.com/2026/03/gtc-2026-recap.html}. Independent analyst notes, open. The NVIDIA roadmap changes of 2026 (Rubin CPX dropped, Kyber shown less dense, Feynman) against which the AMD and Intel roadmaps are read. Slide 32.
\item \textbf{ServeTheHome (2026).} \textit{NVIDIA Vera Rubin NVL72 rack at Hot Chips 2026}. \url{https://www.servethehome.com/nvidia-vera-rubin-nvl72-rack-at-hot-chips-2026/}. Press, secondary source. The rack as presented in the same Hot Chips session as AMD's MI400 talks: NVLink 6 switch with in-network compute, 800 V DC, cable-free trays. Slides 22 and 31.
\item \textbf{Wikipedia (2026).} \textit{AMD Instinct}. \url{https://en.wikipedia.org/wiki/AMD_Instinct}. Tertiary source. A dated pointer page for the Instinct generations; useful for finding primary sources, not for quoting. Slide 31.
\end{itemize}
